\begin{figure}[t]
\center
\includegraphics[width=.16\linewidth]{frostbite_1.png}%
\includegraphics[width=.16\linewidth]{frostbite_2.png}%
\includegraphics[width=.16\linewidth]{frostbite_3.png}%
\hfill
\includegraphics[width=.16\linewidth]{krull_1.png}%
\includegraphics[width=.16\linewidth]{krull_2.png}%
\includegraphics[width=.16\linewidth]{krull_3.png}
\caption{Three consecutive levels in the games \textit{Frostbite} (left) and \textit{Krull} (right). In our experiments, the world model struggles to simulate subsequent levels in \textit{Frostbite}, but not in \textit{Krull}. Indeed, exiting the first level in \textit{Frostbite} requires a long and unlikely sequence of actions to first build the igloo, and then go back to it from the bottom of the screen. Such rare events prevent the world model from internalizing new aspects of the game, which will therefore not be experienced by the policy in imagination. While \textit{Krull} features more diverse levels, the world model successfully reflects this variety, and \textsc{iris} even sets a new state of the art in this environment. This is likely due to more frequent transitions from one stage to the next in \textit{Krull}, resulting in a sufficient coverage of each level.}
\vspace{0.0\linewidth}
\label{fig:frostbite_krull}
\end{figure}
